% Ad Atticum 16.7 (ad-atticum-16-07) --- English translation
% Translated by Claude (Anthropic) from the Latin (Perseus).
% Style guide: STYLE.md.

\ciceroLetterOpener{CICERO ATTICO salutem}

\ciceroSection{1} On the sixth of August, when I had set out from Leucopetra (for from there I was crossing over) and had made some three hundred stadia, I was driven back by a violent south wind to the same Leucopetra. There, while I was waiting for a wind --- for it was the villa of our friend Valerius, where I could be at ease and on familiar terms --- some distinguished men of Regium arrived, fresh from Rome, and among them a host of our friend Brutus, who had left Brutus at Naples. They brought this news: the edict of Brutus and Cassius; that the Senate would meet in full session on the Kalends; that Brutus and Cassius had sent letters to the consulars and praetorians asking them to attend. They reported the highest hope that Antony would give way, that affairs would be settled, that our side would return to Rome. They added too that I was missed --- that I was being half-blamed for it. When I had heard this, without a moment's hesitation I threw aside my plan of departure --- a plan, by Hercules, in which I had never before taken any pleasure. When I read your letter, I was indeed astonished that you had so vehemently changed your opinion, but I judged it not to be without reason. Even if you had not been the urger and prompter of my departure, you had at any rate approved it --- provided I should be at Rome on the Kalends of January. So it was working out that, while there seemed less danger, I would be away; and into the very flame I would be coming. But these things, even if not prudent, are still \textit{[Greek: anemesēta]} --- past blame --- in the first place because they were done on my own opinion, and then, even if on your authority, what is one who gives counsel bound to deliver beyond good faith?

\ciceroSection{3} What I could not sufficiently wonder at was what you wrote in these words: ``Well done then, you who \textit{[Greek: euthanasian]} --- a good ending --- well done! Abandon your country.'' Was I abandoning her? --- or did I seem to you to be abandoning her then? You not only did not forbid me to do so, you actually approved it. What follows is graver still: you say I should polish up for you some \textit{[Greek: scholion]} --- some little dissertation --- on the propriety of my doing what I did. Really, my Atticus? Does my act need a defence --- and that before you, who approved it with such admiration? I shall, to be sure, draw up that \textit{[Greek: apologismon]} --- that justification --- but for one of those against whose unwilling protest I set out. And yet, what need now is there of any \textit{[Greek: scholion]}? If I had persisted, there would have been need. But this very turning back: it is not inconstancy. No learned man (and much has been written on this matter) has ever called a change of plan inconstancy.

\ciceroSection{4} Next, then, comes this: for if you were a pupil of our friend Phaedrus, the excuse would be ready to hand; but as it is, what answer do we give? So, was that my act which I could not justify to Cato? Full of scandal, of course, and disgrace! If only it had seemed so to you from the first! You would have been to me what you usually are --- you would have been my Cato.

\ciceroSection{5} The last point is the most painful. ``For our friend Brutus is silent'' --- that is, he does not dare to admonish a man of my years. I have nothing else that I can suppose to be meant by these words of yours, and by Hercules it is the truth. For on the seventeenth of August, when I had come to Velia, Brutus heard of it --- for he was with his own ships at the river Heles, three miles short of Velia. On foot he came to me at once. Immortal gods, how warmly he rejoiced at my return --- or rather my turning back --- and poured out all that he had kept silent before, so that I remembered that word of yours: ``For our friend Brutus is silent.'' What chiefly grieved him was that I had not been in the Senate on the Kalends of August. Piso he praised to the skies; and as for himself, he rejoiced that I had escaped two very great reproaches --- one, which I knew I was incurring by making the journey, that of despair and of throwing over the state (for crowds wept with me and complained, men whom I could not satisfy with the promise of my swift return); the other, of which Brutus and those with him (and there were many) rejoiced that I had escaped: the reproach of being thought to be travelling to Olympia. For there is nothing more disgraceful, at any season of the state, than this --- and at this season it is \textit{[Greek: anapologēton]} --- past all defence. I myself give wondrous thanks to the south wind, which turned me back from so great a dishonour.

\ciceroSection{6} These you have as the plausible reasons for my turning back --- just ones, indeed, and weighty. But none is juster than this, that you yourself, in another letter, said: ``See to it, if anything is owed to anyone, that there is wherewith to settle like with like.'' For a remarkable \textit{[Greek: dyschrēstia]} --- an awkwardness of credit --- has set in on account of the fear of arms. I read this letter in mid-channel, and what I could see to about it would not come into my head, save that I must defend myself in person and on the spot. But enough of this; the rest face to face.

\ciceroSection{7} I have read Antony's edict, given me by Brutus, and the splendid counter-script of these men; but what those edicts come to, or what they aim at, I do not plainly see. Nor do I now, as Brutus thought, come there to take the state in hand. For what can be done? Did anyone vote with Piso? Did Piso himself return the next day? But they say a man of my age ought not to be far from the grave.

\ciceroSection{8} But I beg you --- what is it that I heard about Brutus? He said you had written that Pilia was \textit{[Greek: peirazesthai paralysei]} --- being assailed by a stroke. I am greatly disturbed. Though you wrote also that you hope for better. So I do, with all my heart; and give her my warmest greetings, and to sweetest Attica too. I wrote this while sailing as I was nearing the Pompeianum, on the nineteenth of August.
